% !PRINTED 19
\weberpage{19}
But this follows immediately from

\weq{7}{Q_1 a^{-1} Q_1 b Q_1 a = Q_1 a^{-1} b a,}

because together with $b$, $a^{-1} b a$ also occurs in $Q$ (because $Q$ is a
normal divisor of $P$). And that the determination of the
indices is also correct is shown by the enumeration. For the degrees of
$P/Q_1$ and $Q/Q_1$ are $\nu \nu_1$ and $\nu_1$, hence the index is $\nu$.

Alongside this there is the following theorem:

\gesperrt{2. If $Q_1$ is a normal divisor of $P$ of degree $q_1$,
and the group $P/Q_1$ itself has a divisor $M$
of degree $m$, then $P$ has a divisor $Q$ of
degree $q = m q_1$. At the same time $Q_1$ is a normal divisor
 of $Q$, and $M = Q/Q_1$. If $M$ is a normal divisor
 of $P/Q_1$, then $Q$ is also a normal divisor
of $P$.}

Let $\mu$ be the index of $Q_1$ with respect to $P$, and let $P$ be
decomposed into the subsidiary groups:

\weq{8}{P = Q_1 + Q_1 e_1 + Q_1 e_2 \cdots + Q_1 e_{\mu - 1},}

so that $e_1, e_2 \ldots e_{\mu - 1}$ are suitably chosen elements in $P$ and

\weq{9}{Q_1,\, Q_1 e_1,\, Q_1 e_2 \ldots Q_1 e_{\mu - 1}}

are the elements of $P/Q_1$. Now if a divisor $M$ is contained in
the group $P/Q_1$, it must consist of a part
of the elements (9), say of

\weq{10}{Q_1,\, Q_1 b_1,\, Q_1 b_2 \ldots Q_1 b_{\nu_1 - 1},}

and if this system is to form a group, then for
any two elements $b_i, b_k$ the product

\weq{11}{Q_1 b_i Q_1 b_k = Q_1 b_i b_k}

must again be contained in (10), hence say $= Q_1 b_h$. If, then,
the system $1, b_1, b_2 \ldots b_{\nu_1 - 1}$ is denoted by $B$, the relation
 (11) can be written, by the method of the composition of the parts,
thus:

\weq{12}{Q_1 B Q_1 B = Q_1 B,}

and so, by \S.\,4, 1., it says that the elements of $P$ contained in

\weq{13}{Q = Q_1 B}

form a group which contains the group
$Q_1$ as a divisor and is itself a divisor of $P$. That
$Q_1$ is a normal divisor of $Q$ follows from \S.\,4, 5., because $Q_1$, as a normal
%CONT>>
